% ---------------------------------------------------------------------------------------
% Diseño de respaldo: solo se usa si MastersDoctoralThesis.cls no está en el proyecto.
% Imita la apariencia de la plantilla Masters/Doctoral Thesis (Palatino, dos caras,
% títulos de capítulo del estilo book, encabezado con línea, páginas de resumen y
% agradecimientos centradas) para que el documento se vea igual sin la clase.
% ---------------------------------------------------------------------------------------
\usepackage[includehead]{geometry}
\usepackage{fancyhdr}
\usepackage[english,spanish,es-noshorthands,es-nodecimaldot,es-nolayout,es-tabla,es-nosectiondot]{babel}
\usepackage{setspace}
\usepackage{longtable}
\usepackage[hidelinks]{hyperref}

% ---- comandos de la plantilla
\newcommand{\thesistitle}[1]{\def\ttitle{#1}}
\newcommand{\supervisor}[1]{\def\supname{#1}}
\newcommand{\examiner}[1]{\def\examname{#1}}
\newcommand{\degree}[1]{\def\degreename{#1}}
\renewcommand{\author}[1]{\def\authorname{#1}}
\newcommand{\addresses}[1]{\def\addressname{#1}}
\newcommand{\subject}[1]{\def\subjectname{#1}}
\newcommand{\keywords}[1]{\def\keywordnames{#1}}
\newcommand{\university}[1]{\def\univname{#1}}
\newcommand{\department}[1]{\def\deptname{#1}}
\newcommand{\group}[1]{\def\groupname{#1}}
\newcommand{\faculty}[1]{\def\facname{#1}}
\newcommand{\HRule}{\rule{\linewidth}{.4mm}}
\newcommand{\addchaptertocentry}[1]{\addcontentsline{toc}{chapter}{#1}}

% ---- resumen y agradecimientos con el diseño centrado de la plantilla
\newcommand{\abstractname}{Resumen}
\newenvironment{abstract}{%
  \cleardoublepage\thispagestyle{plain}
  \begin{center}
  {\normalsize\scshape\univname\par}\bigskip
  {\huge\textit{\abstractname}\par}\bigskip
  {\normalsize\facname\par}{\normalsize\deptname\par}\bigskip
  {\normalsize\degreename\par}\bigskip
  {\normalsize\bfseries\ttitle\par}\medskip
  {\normalsize por \authorname\par}\bigskip
  \end{center}}{\vfil\null}
\newcommand{\acknowledgementname}{Agradecimientos}
\newenvironment{acknowledgements}{%
  \cleardoublepage\thispagestyle{plain}
  \begin{center}{\huge\textit{\acknowledgementname}\par}\end{center}\bigskip}{\vfil\null}
\newcommand{\abbrevname}{Lista de abreviaturas}
\newenvironment{abbreviations}[1]{\cleardoublepage\chapter*{\abbrevname}\begin{longtable}{#1}}{\end{longtable}}
\newcommand{\dedicatory}[1]{\cleardoublepage\thispagestyle{plain}\vspace*{.2\textheight}\begin{center}\itshape #1\end{center}}

% ---- encabezados (con línea bajo el encabezado, como la opción headsepline)
\fancypagestyle{thesis}{\fancyhf{}%
  \fancyhead[LE,RO]{\thepage}%
  \fancyhead[RE]{\itshape\nouppercase{\leftmark}}%
  \fancyhead[LO]{\itshape\nouppercase{\rightmark}}%
  \renewcommand{\headrulewidth}{0.4pt}}
\fancypagestyle{plain}{\fancyhf{}\fancyfoot[C]{\thepage}\renewcommand{\headrulewidth}{0pt}}

% ---- títulos de capítulo con el estilo de la plantilla, pero con menos espacio vertical
\makeatletter
\def\@makechapterhead#1{\vspace*{10\p@}{\parindent\z@\raggedright\normalfont
  \ifnum\c@secnumdepth>\m@ne\if@mainmatter\huge\bfseries\@chapapp\space\thechapter\par\nobreak\vskip 12\p@\fi\fi
  \interlinepenalty\@M\Huge\bfseries #1\par\nobreak\vskip 24\p@}}
\def\@makeschapterhead#1{\vspace*{10\p@}{\parindent\z@\raggedright\normalfont
  \interlinepenalty\@M\Huge\bfseries #1\par\nobreak\vskip 24\p@}}
\makeatother
